% Two-stage factorisation of the CY-to-chiral functor $\Phi_d$.
% Structural refinement surfaced between Wave 15 and Wave 16.
% Author: Raeez Lorgat.
% Chriss--Ginzburg register. Subscripted $\kappa_\bullet$. No AI attribution.

\providecommand{\ClaimStatusTheorem}{\textsc{[Theorem]}}
\providecommand{\ClaimStatusConjectured}{\textsc{[Conjectured]}}
\providecommand{\ClaimStatusCorrected}{\textsc{[Corrected]}}
\providecommand{\ClaimStatusDefinition}{\textsc{[Definition]}}
\providecommand{\kBKM}{\kappa_{\mathrm{BKM}}}
\providecommand{\kch}{\kappa_{\mathrm{ch}}}
\providecommand{\kcat}{\kappa_{\mathrm{cat}}}
\providecommand{\kfib}{\kappa_{\mathrm{fiber}}}
\providecommand{\CYcat}{\mathrm{CY}\text{-cat}}
\providecommand{\ChirAlg}{\mathrm{ChirAlg}}
\providecommand{\HolFA}{\mathrm{HolFA}}
\providecommand{\Sp}{\mathrm{Sp}}
\providecommand{\ZZ}{\mathbb{Z}}
\providecommand{\hCS}{\mathrm{hCS}}
\providecommand{\Obs}{\mathrm{Obs}}
\providecommand{\CE}{\mathrm{CE}}
\providecommand{\cE}{\mathcal{E}}
\providecommand{\cO}{\mathcal{O}}
\providecommand{\bH}{\mathbf{H}}

\section{The two-stage factorisation $\Phi_d = \Phi^{\mathrm{FA}}_d \circ \Sp_{\Sigma_{d-1},C}$}
\label{wn:sec:two-stage-Phi-factorisation}

\emph{The CY-to-chiral functor $\Phi_d$ factors canonically through the holomorphic
$E_d$-factorisation algebra on the CY$_d$ itself. The $E_n$-chiral shadow on a
curve is not the primary output but a specialisation, selected by a choice of
$(d{-}1)$-cycle $\Sigma_{d-1} \subset X$ and reference curve $C \subset X$.
Borcherds, Igusa, and Fake-Monster are sibling $E_1$-specialisations of the same
$E_3$-holomorphic factorisation algebra.}

\subsection{Stage 1: canonical $E_d$-hFA assignment}
\label{wn:subsec:stage-1-Ed-hFA}

\begin{definition}[Stage 1, canonical]\ClaimStatusDefinition
\label{wn:def:stage-1-Ed-hFA}
For a smooth projective Calabi--Yau $d$-fold $X$, the first stage
$$
\Phi^{\mathrm{FA}}_d: \CYcat_d \longrightarrow E_d\text{-}\HolFA(X)
$$
assigns to $X$ (equipped with its CY datum $(\Omega_X, g_X, \ldots)$) the
holomorphic $E_d$-factorisation algebra of Costello--Gwilliam Vol~II
observables of the 6D (resp.\ $2d$-D) holomorphic Chern--Simons theory on $X$,
with gauge data from the CY category. The assignment is canonical up to
contractible choice: the operadic $E_d$-structure is pinned by Kontsevich--Tamarkin
formality of $\mathcal{D}_d$ (Kontsevich 1999, Tamarkin 2003), and the
factorisation/locality axioms by Costello--Gwilliam (2017, Vol~II \S 3)
with Costello--Li (2016 \emph{JHEP}, hCS specialisation) supplying the
explicit $\bar\partial$-twist at $d=3$.
\end{definition}

\begin{theorem}[Stage 1 is canonical]\ClaimStatusTheorem
\label{wn:thm:stage-1-canonical}
Stage 1 is determined up to contractible choice by:
(i) Kontsevich 1999 / Tamarkin 2003 $E_d$-formality of the Dolbeault little-$d$-disks operad;
(ii) Costello--Gwilliam 2017 Vol~II Thm.~9.3.1 (factorisation $\leftrightarrow$ operadic $E_d$);
(iii) Costello--Li 2016 holomorphic twist / BCOV realisation at $d=3$;
(iv) Francis--Gaitsgory 2011 Koszul duality for $E_d$-algebras
(for the intrinsic coalgebra description);
(v) Gwilliam--Williams 2021 Thm.~2.5.5 $E_d^{\mathrm{hol}}$-to-$E_d$ comparison.
\end{theorem}

\begin{remark}[What Stage 1 does \emph{not} involve]
\label{wn:rmk:stage-1-no-curve-choice}
The output $\Phi^{\mathrm{FA}}_d(X)$ lives on $X$ itself, not on any curve
$C \subset X$. No $(d{-}1)$-cycle is chosen. No reference disk is fixed. No
specialisation occurs. The 6D hCS observable $E_3$-algebra established in
Wave~13~E1--E5 is Stage 1 at $d=3$ for the specific CY category of
$\mathbb{C}^3$ (open) or $K3 \times E$ (compact).
\end{remark}

\subsection{Stage 2: specialisation via factorisation homology}
\label{wn:subsec:stage-2-specialisation}

\begin{definition}[Stage 2, specialisation]\ClaimStatusDefinition
\label{wn:def:stage-2-specialisation}
Given a closed $(d{-}1)$-dimensional submanifold (or subvariety)
$\Sigma_{d-1} \subset X$ and a reference curve $C \subset X$,
$$
\Sp_{\Sigma_{d-1},C}: E_d\text{-}\HolFA(X) \longrightarrow E_1\text{-}\ChirAlg(C)
$$
is factorisation homology over $\Sigma_{d-1}$ fibrewise along the projection
$X \to C$ (or, equivalently, $\int_{\Sigma_{d-1}}$ evaluated on the slice
over $C$). Concretely, for an $E_d$-hFA $\mathcal{F}$ on $X$,
$$
\Sp_{\Sigma_{d-1},C}(\mathcal{F})(U) := \int_{\Sigma_{d-1}} \mathcal{F}\big|_{U \times \Sigma_{d-1}}
$$
for $U \subset C$ open. The output is an $E_1$-chiral algebra on $C$ via
Lurie \emph{HA} §5.3 Dunn additivity ($E_d = E_{d-1} \otimes E_1$ after
fibrewise integration).
\end{definition}

\begin{definition}[$\Phi_d$ as composite]\ClaimStatusDefinition
\label{wn:def:Phi-d-composite}
The CY-to-chiral functor decomposes as
$$
\Phi_d = \Sp_{\Sigma_{d-1},C} \circ \Phi^{\mathrm{FA}}_d :
\CYcat_d \longrightarrow E_1\text{-}\ChirAlg(C).
$$
Stage 1 is pinned up to contractible choice; Stage 2 depends on a choice of
$(\Sigma_{d-1}, C) \subset X$. The assignment $X \mapsto \mathcal{F}_X$ is
the canonical intrinsic object; the chiral-algebra-on-a-curve shadow
$\Sp_{\Sigma_{d-1},C}(\mathcal{F}_X)$ is one of multiple specialisations.
\end{definition}

\subsection{Per-dimension dictionary}
\label{wn:subsec:per-dim-dictionary}

\begin{center}
\renewcommand{\arraystretch}{1.22}
\begin{tabular}{|c|l|l|l|}
\hline
$d$ & Native $E_d$-hFA on CY$_d$ & Specialisation cycle $\Sigma_{d-1}$ & $E_1$-chiral shadow \\
\hline
$1$ & $E_1$-hFA on $E_\tau$ (elliptic curve) & trivial $\Sigma_0 = \mathrm{pt}$ & $E_\infty$ Heisenberg \\
$2$ & $E_2$-hFA on K3 (surface factorisation) & $S^1 \subset K3$ (Nakajima cycle) & $E_2$-braided Mukai $H_{\mathrm{Muk}}$ \\
$3$ & $E_3$-hFA on $K3\times E$ (Costello--Li hCS) & $K3$ fibre & $E_1$-ordered chiral on $E$ (CY-A$_3$) \\
$3$ & $E_3$-hFA on $A \times E$ (abelian$\times$E) & $T^4$ fibre & distinct $E_1$-ordered chiral \\
$5$ & $E_5$-hFA on $K3 \times K3 \times E$ & $K3 \times K3$ & $E_1$-ordered chiral on $E$ \\
\hline
\end{tabular}
\end{center}

\subsection{Consequence: ``many BKMs from one CY$_3$'' is a theorem}
\label{wn:subsec:many-BKMs-theorem-grade}

\begin{conjecture}[Sibling specialisation]\ClaimStatusConjectured
\label{wn:conj:sibling-specialisation}
The Borcherds monster denominator algebra, the Igusa $\Phi_{10}$ GBKM
$\mathfrak{g}_{\Delta_5}$, and the Fake Monster Lie algebra of signature
$\Lambda^{1,25}$ are sibling $E_1$-specialisations of a common $E_3$-hFA
framework: each arises as $\Sp_{\Sigma_2, C}(\mathcal{F}_{X_i})$ for an
appropriate CY$_3$ datum $X_i$ and $(\Sigma_2, C)$ choice. The multiplicity of
$E_1$-specialisations is a feature encoding the
$H^0(X, \mathcal{M}_{\Sigma_2} \otimes \mathcal{M}_C)$ moduli parameters of
the specialisation datum, not a defect of the framework.
\end{conjecture}

\begin{remark}[Relation to the $\{2,3,5,24\}$ quadruple on $K3\times E$]
\label{wn:rmk:2-3-5-24-reinterpretation}
The four values $\{2, 3, 5, 24\}$ for $K3\times E$ are
\emph{not} four $\Sigma_2$-specialisations of one $E_3$-hFA. They are four
distinct constructions measuring four distinct invariants: $\kcat$ (Mukai
lattice rank--22 complement), $\kch$ (Hodge supertrace identification),
$\kBKM$ (Borcherds weight $c_1(0)/2$ of $\Delta_5$), $\kfib$ (fibre rank of
the K3 lattice). The two-stage factorisation refines the $\kch$ row:
$\kch(K3\times E) = 0$ is the invariant of
$\mathcal{F}_{K3\times E} \in E_3$-$\HolFA(K3\times E)$, independent of the
specialisation $(\Sigma_2, C)$. Different $\Sigma_2$ choices do not change
$\kch$; they change which $E_1$-chiral-algebra-on-a-curve shadow is
visible, which is a richer structure than a single numeric invariant.
\end{remark}

\subsection{Programme-level rectification protocol}
\label{wn:subsec:programme-rectification}

\begin{itemize}
\item \textsc{Rectification~A} \ClaimStatusCorrected: every prior claim of the
form ``$\Phi_d$ natively produces an $E_n$-chiral algebra on a curve at
$d = k$'' is rewritten: ``$\Phi_d$ natively produces an $E_d$-hFA on the
CY$_d$; a choice of $\Sigma_{d-1}$-cycle and reference curve $C$ specialises
this to an $E_1$-chiral algebra on $C$.''

\item \textsc{Rectification~B} \ClaimStatusCorrected: every invocation of
``the ambient $E_n$-chiral algebra'' with $n = d$ ($d \geq 2$) is
re-sourced: the ambient structure is the $E_d$-hFA on $X$; the $E_n$-chiral
algebra appears on the curve after Stage~2.

\item \textsc{Rectification~C} \ClaimStatusConjectured: the seven-face
$r_{\mathrm{CY}}$ picture on $K3 \times E$
is reinterpreted
as seven $(\Sigma_2, C)$-specialisation data of the same
$\mathcal{F}_{K3\times E}$, not seven distinct $\Phi$ applications. The
seven $r_{\mathrm{CY}}$ faces are distinguished by cycle choice, not by
functorial ambiguity.

\item \textsc{Rectification~D} \ClaimStatusConjectured: the universal
Borcherds weight identity $\kBKM(\Phi_N) = c_N(0)/2$ on
$N \in \{1,2,3,4,6\}$ is rewritten as a statement about the specialisation
invariant
$$
\kBKM\!\left(\Sp_{\Sigma_2, C}(\mathcal{F}_{X_N})\right) = c_N(0)/2
$$
for the canonical $(\Sigma_2, C) = (K3, E)$ choice on the CHL quotient
$X_N = (K3 \times E)/(\ZZ/N\ZZ)$.
\end{itemize}

\subsection{Connection to Costello--Gwilliam--Li and Stage 1 primary sources}
\label{wn:subsec:stage-1-primary-sources}

Stage 1 consumes the full Costello machine. The relevant primary references:
\begin{enumerate}
\item \textbf{Costello 2011} \emph{Renormalization and Effective Field Theory}:
existence of BV-quantised observables as factorisation algebras.
\item \textbf{Costello--Gwilliam 2017} \emph{Factorization Algebras in
Quantum Field Theory} Vol~I (classical BV, factorisation axioms),
Vol~II (quantum, hCS on $\mathbb{C}^3$, operadic $E_d$-comparison).
\item \textbf{Costello--Li 2016} ``Quantization of open-closed BCOV theory
I'' \emph{JHEP}: hCS on CY$_3$ as the holomorphic twist of
$\mathcal{N} = 2$ on $X \times \mathbb{R}^4$.
\item \textbf{Costello 2013} arXiv:1303.2632: 4D CS (shadow at $d=1$ via
$\Sp_{\mathrm{pt}}$) with Yangian quantisation.
\item \textbf{Costello--Dimofte--Paquette 2020} arXiv:2005.00083: 5D hCS
boundary chiral-algebra-of-operators, which is a Stage-2 shadow of the
6D hCS $E_3$-hFA under boundary pinning.
\item \textbf{Gwilliam--Williams 2021} arXiv:2009.05037:
$E_d^{\mathrm{hol}}$ vs.\ $E_d$ comparison theorem securing Stage~1's
operadic precision.
\end{enumerate}

\subsection{Next mathematical target}
\label{wn:subsec:two-stage-next-target}

The frontier opened by the two-stage factorisation: for each
$(\Sigma_{d-1}, C)$-pair on $K3 \times E$, compute the chiral-algebra shadow
$\Sp_{\Sigma_2, C}(\mathcal{F}_{K3\times E})$ and verify it lands in a
known $E_1$-chiral algebra. Partial data:
$(\Sigma_2 = K3, C = E) \mapsto \bH_{\Delta_5}$ chiral bialgebra (Wave~11
meta-frontier). Open: $(\Sigma_2, C)$ choices outside the canonical
$(K3, E)$ --- does $(\Sigma_2 = $ elliptic surface, $C = $ base$\,\mathbb{P}^1)$
yield a new $E_1$-chiral, and is it Gritsenko--Cléry-indexed? This
connects to Wave~15~M3's $\Lambda^{3,2} \hookrightarrow \Lambda^{3,3}$
envelope: the larger envelope GBKM $\mathfrak{g}^{\mathrm{BKM}}_{\Lambda^{3,3}}$
emerges from a non-standard $\Sigma_2$-choice.
